\documentclass[main.tex]{subfiles}


\begin{document}

% \AddToShipoutPictureFG{%
% \begin{tikzpicture}[overlay,remember picture]
% \path (current page.south west) -- (current page.north east)
% node[midway,scale=1,color=gray,sloped,opacity=0.4] {\textnormal{Кравченко Игорь Игоревич\hspace{1cm} +79010144910\hspace{1cm} super.antig@yandex.ru}};
% \end{tikzpicture}
% }


% %from pagenumber to up
% \phantomsection 
% \hypertarget{toc}{}

 

 

% номер секции вручную
\setcounter{section}{10
}
\addtocounter{section}{-1} 

\section{Свободное падение}


\sbox{0}{%
    \begin{tikzpicture}[font=\small,            line cap=round,                             >={Stealth[inset=0pt,round,length=3mm, angle'=20]},vec/.style={>={Stealth[inset=0pt,round,length=3.5mm, angle'=20]}}]


\filldraw[lightgray,semitransparent] (5,0) rectangle (7,-0.3);
\draw (5,0) -- (7,0);

\shade[ball color=lightgray,nearly transparent,yshift=4cm] (6,.3) circle (.3);
\shade[ball color=lightgray,semitransparent,yshift=3cm] (6,.3) circle (.3);
\shade[ball color=lightgray] (6,.4) circle (.3);

%vec
\node at (6,3.3) [circle,fill=NavyBlue,inner sep=0pt,minimum size=1mm,label=below:] {};
\draw[->,NavyBlue,thick,vec] (6,3.3) --++ (0,-.75) node[black,right] {$\vec v_1$};

\node at (6,0.4) [circle,fill=NavyBlue,inner sep=0pt,minimum size=1mm,label=below:] {};
\draw[->,NavyBlue,thick,vec] (6,0.4) --++ (0,-1.5) node[black,right] {$\vec v_2$};


    \end{tikzpicture}%
}

{%
\setlength\intextsep{0pt}
\begin{wrapfigure}[17]{r}{\wd0}
    \centering
    \usebox{0}%
    \caption{Падение шара}
    \label{fig:p19}
\end{wrapfigure}

\textbf{Свободным падением} называют движение, происходящее под влиянием только планеты. Влияние планеты выражается, так сказать, в стремлении всех тел падать вниз. Происходит такое движение обычно вблизи земной поверхности~--- то есть поверхность считается плоской.  

Опыт с падением тяжелого шара на поверхность планеты показан на рис.\,\ref{fig:p19} (положение тела фиксируют ежесекундно).
    
Из рис.\,\ref{fig:p19} видно, что за каждую секунду движения скорость тела меняется на одну и ту же величину, то есть движение является равноускоренным. \textbf{Ускорение свободного падения} обозначается~$\vec g$, оно напралено вертикально вниз (или к центру планеты).

В большинстве задач вблизи поверхности планеты ускорение~$\vec g$ для любого тела полагают одинаковым независимо от, например, его размеров, массы и т.\,д.

Теперь можно рассмотреть полет тела, пущенного с начальной скоростью под некоторым углом к горизонту (рис.\,\ref{fig:p20}).

}

\begin{figure}[H]
\centering

    \begin{tikzpicture}[font=\small,            line cap=round,                             >={Stealth[inset=0pt,round,length=3mm, angle'=20]},vec/.style={>={Stealth[inset=0pt,round,length=3.5mm, angle'=20]}}]

\def\initialV{6} % initial speed
\def\launchangle{60} % angle of launch
\def\gravity{10} % force of gravity
\def\yaxemax{4}
\def\xaxemax{6}
\coordinate (V) at ({1+\initialV*cos(\launchangle)},{1+\initialV*sin(\launchangle)});
\coordinate (Vx) at ({1+\initialV*cos(\launchangle)},1);

\coordinate [label= below left :{$O$}] (O) at (0,0);
\coordinate (init) at (1,1) node [below left] {};

%axes and patterns
\filldraw[lightgray,semitransparent] (0,0) rectangle (5,-0.3);
\draw[->] (0,0) -- node[left,at end] {$y$} (0,\yaxemax);
\draw[->] (0,0) -- node[below,at end] {$x$} (\xaxemax,0);

 % the lower plane is one unit below the exit, so set y = t*v_0*sin(theta) - 0.5*g*t^2 = -1, and solve for t
\pgfmathsetmacro\timeofimpact{\initialV*sin(\launchangle)/\gravity + sqrt((\initialV*sin(\launchangle)/\gravity)^2 + 2/\gravity)}

%t for ymax of plot
\pgfmathsetmacro\tforymax{-\initialV*sin(\launchangle)/-\gravity}


%points
\coordinate [label= left :{$y_{\max}$}] (ymax) at (0,{1+\tforymax*\initialV*sin(\launchangle) - 0.5*\gravity*\tforymax^2});
\coordinate [label= below :{}] (pointmax) at ({1+\tforymax*\initialV*cos(\launchangle)},{1+\tforymax*\initialV*sin(\launchangle) - 0.5*\gravity*\tforymax^2});

%dashed
\draw[-,densely dashed] (ymax) -- (pointmax);
\draw[-,densely dashed] (0,1) node [left] {$y_0$} -- ({1+\tforymax*\initialV*cos(\launchangle)},1);
\draw[-,densely dashed] (1,1) -- (1,0) node [above right] {$x_0$};



 % domain based on calculation above
\draw [thick,samples=100,variable=\t,domain=0:\timeofimpact] (1,1) plot ({1+\t*\initialV*cos(\launchangle)},{1+\t*\initialV*sin(\launchangle) - 0.5*\gravity*\t^2}) node [above right] {$x_{\max}$};

%v
\draw[->,NavyBlue,thick,vec] (1,1) --++ ({\launchangle}:2.5) node[left=1mm,black] {$\vec v_0$};
\pic [draw,angle radius=10mm,angle eccentricity=1.25,"$\alpha$"] {angle = Vx--init--V};

%g
\draw[->,red,thick,vec] ({1+\timeofimpact*\initialV*cos(\launchangle)},\yaxemax) --++ ({-90}:2) node[right,midway,black] {$\vec g$};

\shade[ball color=gray] (1,1) circle (1mm);


    \end{tikzpicture}
\caption{Полет тела}
\label{fig:p20}
\end{figure}

В начальный момент времени тело находилось над поверхностью планеты в точке с координатами~$x_0$ и~$y_0$ и имело скорость~$\vec v_0$, направленную под углом~$\alpha$ к горизонту. Ускорение свободное падения~$\vec g$ есть красный вектор на рис.\,\ref{fig:p20}. Можно показать, что траекторией тела в системе координат~$xOy$ является парабола (черная жирная линия на рис.\,\ref{fig:p20}). Конечной точкой падения будет точка пересечения траектории с планетной поверхностью~--- то есть с осью~$Ox$. Полезно замечать характерные точки <<следа>> тела: так, точка наивысшего подъема характеризуется координатой~$y_{\max}$; а точка наибольшего удаления по горизонтали~--- координатой~$x_{\max}$. 

Вот опорные уравнения для построения решений многих задач этой темы. 
\begin{gather*}
\left\{
\begin{aligned}
&x=x_0+v_{0x}t+\dfrac{g_xt^2}{2},\\
&y=y_0+v_{0y}t+\dfrac{g_yt^2}{2}.
\end{aligned}
\right.
\\
\left\{
\begin{aligned}
&v_x=v_{0x}+g_xt,\\
&v_y=v_{0y}+g_yt.
\end{aligned}
\right.
\end{gather*}





 
\end{document}